\section*{Chapitre \ref{ch:g2:water-on-earth} --- L'eau sur la Terre}

\begin{solution}{exo:g2:water-on-earth:1}
\omterm{def:g2:water-on-earth:salt-fresh}{Eau salée}~: la \omterm{def:g2:water-on-earth:ocean}{mer}, l'\omterm{def:g2:water-on-earth:ocean}{océan}. \omterm{def:g2:water-on-earth:salt-fresh}{Eau douce}~: la \omterm{def:g1:landscapes-around-us:water}{rivière}, la pluie, le \omterm{def:g1:landscapes-around-us:water}{lac} de
\omterm{def:g1:landscapes-around-us:landform}{montagne}, la glace sur la \omterm{def:g1:landscapes-around-us:landform}{montagne}.
\end{solution}

\begin{solution}{exo:g2:water-on-earth:2}
Parce que la plus grande partie de sa surface est couverte par l'eau bleue
des \omterm{def:g2:water-on-earth:ocean}{océans}~: vue de l'espace, elle paraît bleue.
\end{solution}

\begin{solution}{exo:g2:water-on-earth:3}
$100 - 71 = 29$ carrés sont des terres.
\end{solution}

\begin{solution}{exo:g2:water-on-earth:4}
Un \omterm{def:g2:water-on-earth:ocean}{océan} est plus grand~: une \omterm{def:g2:water-on-earth:ocean}{mer} est une étendue d'\omterm{def:g2:water-on-earth:salt-fresh}{eau salée} plus petite,
souvent au bord d'un \omterm{def:g2:water-on-earth:ocean}{océan}.
\end{solution}

\begin{solution}{exo:g2:water-on-earth:5}
Non~: l'eau de l'\omterm{def:g2:water-on-earth:ocean}{océan} est de l'\omterm{def:g2:water-on-earth:salt-fresh}{eau salée}, et la boire fait perdre de l'eau
au corps. Le marin doit emporter de l'\omterm{def:g2:water-on-earth:salt-fresh}{eau douce}.
\end{solution}

\begin{solution}{exo:g2:water-on-earth:6}
Par exemple~: dans les \omterm{def:g1:landscapes-around-us:water}{rivières}, dans les \omterm{def:g1:landscapes-around-us:water}{lacs}, dans la glace et la neige
des \omterm{def:g1:landscapes-around-us:landform}{montagnes} et de l'Antarctique, sous la terre, dans la pluie.
\end{solution}

\begin{solution}{exo:g2:water-on-earth:7}
\omterm{def:g2:water-on-earth:salt-fresh}{Eau salée}~: 96 seaux. \omterm{def:g2:water-on-earth:salt-fresh}{Eau douce}~: $2 + 2 = 4$ seaux.
(Pour être exact, une partie de l'eau sous la terre est salée elle aussi~;
l'image la compte avec l'\omterm{def:g2:water-on-earth:salt-fresh}{eau douce}.)
\end{solution}

\begin{solution}{exo:g2:water-on-earth:8}
Les \omterm{def:g2:water-on-earth:ocean}{océans} (96 seaux), puis la glace et la neige (2 seaux), puis les
\omterm{def:g1:landscapes-around-us:water}{rivières} (bien moins d'un seau).
\end{solution}

\begin{solution}{exo:g2:water-on-earth:9}
Elle peut s'infiltrer dans le sol~; elle peut ruisseler sur le sol jusqu'à
un ruisseau ou une \omterm{def:g1:landscapes-around-us:water}{rivière}~; elle peut s'évaporer et retourner dans l'air
quand le soleil revient.
\end{solution}

\begin{solution}{exo:g2:water-on-earth:10}
Il reste du sel au fond de la casserole. L'eau s'est évaporée~: elle est
partie dans l'air sous forme de vapeur d'eau.
\end{solution}

\begin{solution}{exo:g2:water-on-earth:11}
Par exemple~: la goutte coule le long du toit jusque dans la gouttière et
l'égout, puis dans un ruisseau et une \omterm{def:g1:landscapes-around-us:water}{rivière}~; la \omterm{def:g1:landscapes-around-us:water}{rivière} l'emporte jusqu'à
la \omterm{def:g2:water-on-earth:ocean}{mer}~; un jour de soleil, elle s'évapore de la \omterm{def:g2:water-on-earth:ocean}{mer}, monte dans l'air et
devient une partie d'un nuage.
\end{solution}
